\section{Piecewise $\mathcal{C}^2$ certificates in dimension one}

The setup is pretty similiar, with the exception that $V$ is now piecewise $\mathcal{C}^2$ instead of piecewise linear. 

\begin{setup}
Let:
\begin{enumerate}
    \item $dX_t = f(X_t) dt + g(X_t) dW_t$ be a one-dimensional SDE.
    \item $V(x)$ be a piecewise $\mathcal{C}^2$ function such that $V(x) \in \mathcal{C}^2$, where $x \in \left( x_i, x_{i+1} \right)$.
    \item $V(x)$ is continuous at breakpoints: $V(x_i^+)=V(x_i^-) = V(x_i)$.
\end{enumerate}

\end{setup}

We will again consider what conditions on $V$ are sufficient to guarantee that $V(X_t)$ is a supermartingale. Again, the natural hypothesis is:

\begin{hypothesis}
Does concavity of $V(x)$ (along with an interior / drift term condition) guarantee that $V(X_t)$ is a supermartingale?
\end{hypothesis}

We apply the Itô-Tanaka formula to $V(X_t)$, and we get similiar terms, with the addition of
\begin{equation}
  C_t = \frac{1}{2} \int_{0}^{t} V''(X_s) g^2(X_s) ds
\end{equation}

Since the integral in the $C_t$ term is over the same intervals as the $D_t$ term, we can combine them to recover the generator / Lie derivative term:
\begin{equation}
  D_t + C_t = \int_{0}^{t} \left( V'(X_s) f(X_s) + \frac{1}{2} V''(X_s) g^2(X_s) \right) ds = \int_{0}^{t} \mathcal{L} V(X_s) ds
\end{equation}

The proof that $V(X_t)$ is a supermartingale is similiar to the piecewise linear case, but now we need to impose the condition:
\begin{equation}
  \mathcal{L} V(x) = V'(x) f(x) + \frac{1}{2} V''(x) g^2(x) \leq 0
\end{equation}

\newpage

